
\subsubsection{3.1 Hypothesis}

The hypothesis under test (H-E1) states: the coefficient of variation (CV) of linear probe accuracy trajectories can distinguish spurious from core features with AUC $\geq$ 0.75.

The underlying reasoning: spurious features, being simpler and correlating with labels regardless of subgroup, should be learned uniformly across sample subsets (low CV). Core features, relevant to specific subgroups, should emerge differentially (high CV).

\subsubsection{3.2 Feature Extraction}

CLIP ViT-B/16 \citep{radford2021learning} served as the feature extractor. For each Waterbirds image, the 512-dimensional embedding from the final layer was extracted and L2-normalized. Features were cached to disk for reproducibility.

\subsubsection{3.3 CV Measurement Protocol}

For each concept (background and bird type), logistic regression probes were trained with regularization strength C $\in$ {0.001, 0.01, 0.1, 1, 10, 100}. The solver was L-BFGS with a maximum of 1000 iterations.

To measure emergence uniformity, probes were trained on 5 random 20\% subsets of training data (seed = 42). For each concept:

\begin{enumerate}
\item Train probes on each subset across all C values
\item Record accuracy at each (subset, C) combination
\item Compute improvement rate: accuracy at maximum C minus accuracy at minimum C
\item Calculate CV of improvement rates across subsets: CV = std(rates) / mean(rates)
\end{enumerate}

\subsubsection{3.4 Evaluation}

Waterbirds provides ground truth: background is the spurious feature (95\% correlated with label) and bird type is the core feature. Binary classification was evaluated using ROC-AUC, with the gate criterion set at AUC $\geq$ 0.75.

\subsubsection{3.5 Experimental Configuration}

\begin{table}[htbp]
\centering
\resizebox{\columnwidth}{!}{%
\begin{tabular}{ll}
\toprule
Parameter & Value \\
\midrule
Dataset & Waterbirds (train split) \\
Train samples & 4,795 \\
Feature extractor & CLIP ViT-B/16 (frozen) \\
Embedding dimension & 512 \\
Probe & Logistic Regression (L-BFGS) \\
C sweep & [0.001, 0.01, 0.1, 1, 10, 100] \\
Number of subsets & 5 \\
Subset fraction & 20\% \\
Random seed & 42 \\
Gate threshold & AUC $\geq$ 0.75 \\
\bottomrule
\end{tabular}
}
\end{table}

